\subsection{径向对称性}\label{sec:radial-symmetry}

在一篇 1979 年的论文中, Gidas, Ni 和 Nirenberg~\cite{GNN} 建立了某些非线性椭圆方程正解的径向对称性. 这一技巧基于最大值原理. 本节考虑其结果中最简单的情形, 即以下定理.

\begin{Theorem}
\label{thm:radial-symmetry-positive-solution}
设 $f:\mathbb{R}\to\mathbb{R}$ 属于 $C^1$ 类. 在球 $B=B_R(0)\subset\mathbb{R}^n$ 内, 设 $u>0$ 为方程
\begin{equation}\label{eq:radial-symmetry-semilinear-equation}
\Delta u+f(u)=0,
\end{equation}
的解 (属于 $C^2(\overline B)$ 类), 并满足边界条件
\begin{equation}\label{eq:radial-symmetry-zero-boundary}
u=0\quad\text{on }\partial B.
\end{equation}
那么 $u$ 径向对称且严格单调递减, 即
\begin{equation}\label{eq:radial-symmetry-strict-decrease}
\dfrac{\partial u}{\partial r}<0\quad\text{for }0<r<R.
\end{equation}

\end{Theorem}

值得注意的是, 无论 $f$ 为何, 这一结果都成立.

\subsubsection{两个辅助引理}\label{subsec:radial-auxiliary-lemmas}

本节余下部分始终取 $B$ 和 $u$ 如 Theorem~\ref{thm:radial-symmetry-positive-solution} 的假设所述.

\begin{Lemma}
\label{lem:radial-boundary-directional-derivative}
设 $\bm{x}_0\in\partial B\cap\{x_1>0\}$. 那么存在某个 $\delta>0$, 使在 $B\cap\{|\bm{x}-\bm{x}_0|<\delta\}$ 内有 $\partial u/\partial x_1<0$.
\end{Lemma}

\begin{Proof}
为简化记号, 用 $u_1$ 表示 $\partial u/\partial x_1$, 用 $u_{11}$ 表示 $\partial^2u/\partial x_1^2$, 等等.

若引理不成立, 则存在点列 $\bm{x}_j\in B$, 满足 $u_1(\bm{x}_j)\geq0$ 且 $\bm{x}_j\to\bm{x}_0$. 另一方面, 由于 $B$ 内 $u>0$ 且边界上 $u=0$, 必有 $u_1(\bm{x}_0)\leq0$. 因而由连续性, 得到 $u_1(\bm{x}_0)=0$. 注意, 由于在 $\partial B$ 上 $u=0$, 所有切向导数均为零. 由于 $x_1$ 方向不在切平面内, 已证明 $\nabla u(\bm{x}_0)=0$.

断言二阶导数 $u_{11}(\bm{x}_0)$ 也为零. 首先, 由于 $u(\bm{x}_0)=u_1(\bm{x}_0)=0$, 且 $B$ 内 $u>0$, 必有 $u_{11}(\bm{x}_0)\geq0$. 现在假设 $u_{11}(\bm{x}_0)>0$. 那么 $u_{11}$ 在 $\bm{x}_0$ 的某个邻域 $N$ 内也为正. 考虑沿正 $x_1$ 方向连接 $\bm{x}_j$ 与 $\partial B$ 上一点 $\bm{y}_j$ 的直线段 $\Gamma$. 对充分大的 $j$, $\Gamma$ 完全位于 $N$ 内, 因而
\begin{equation}\label{eq:radial-boundary-second-derivative-segment}
u_1(\bm{y}_j)-u_1(\bm{x}_j)=\int_\Gamma u_{11}(\bm{x})\,\mathrm{d}x_1
\end{equation}
为正. 但由于 $u_1(\bm{x}_j)\geq0$ 且 $u_1(\bm{y}_j)\leq0$, 这产生矛盾.

现在假设 $f(0)\geq0$. 则有
\begin{equation}\label{eq:radial-nonlinearity-zero-value-shift}
\Delta u+f(u)-f(0)\leq0,
\end{equation}
由中值定理可找到函数 $c(\bm{x})$, 使 $f(u)-f(0)=c(\bm{x})u$. 将 Lemma~\ref{lem:elliptic-boundary-point} 应用于 $-u$, 得到 $u_1(\bm{x}_0)<0$, 矛盾.

因此假设 $f(0)<0$. 那么 $\Delta u(\bm{x}_0)=-f(0)>0$. 另一方面, 由于在 $\partial B$ 上 $u=0$ 且 $\nabla u(\bm{x}_0)=0$, 容易验证 $u_{11}(\bm{x}_0)=n_1^2\Delta u(\bm{x}_0)\ne0$, 其中 $\bm{n}$ 表示 $\partial B$ 的单位法向 (见 Problem~\ref{ex:boundary-hessian-radial-identity}). 再次产生矛盾.
\end{Proof}

对 $\lambda\in\mathbb{R}$, 设 $T_\lambda$ 表示平面 $x_1=\lambda$, 令 $\Sigma(\lambda)=B\cap\{x_1>\lambda\}$. 此外, 设 $\Sigma^{\prime}(\lambda)$ 表示 $\Sigma(\lambda)$ 关于 $T_\lambda$ 的反射, 设 $\bm{x}^\lambda$ 表示点 $\bm{x}$ 关于 $T_\lambda$ 的反射. 有以下引理.

\begin{Lemma}
\label{lem:radial-reflection-strictness}
假设对某个 $\lambda\in[0,R)$ 有
\begin{equation}\label{eq:radial-reflection-weak-comparison}
u_1(\bm{x})<0,\quad u(\bm{x})\leq u(\bm{x}^\lambda)\quad\forall\bm{x}\in\Sigma(\lambda),
\end{equation}
但在 $\Sigma(\lambda)$ 内 $u(\bm{x})$ 不恒等于 $u(\bm{x}^\lambda)$. 那么在整个 $\Sigma(\lambda)$ 内 $u(\bm{x})<u(\bm{x}^\lambda)$, 且在 $B\cap T_\lambda$ 上 $u_1(\bm{x})<0$.
\end{Lemma}

\begin{Proof}
在 $\Sigma^{\prime}(\lambda)$ 内, 令 $v(\bm{x})=u(\bm{x}^\lambda)$. 那么 $v$ 满足方程 $\Delta v+f(v)=0$. 令 $w=v-u$, 并取 $c(\bm{x})$, 使 $f(v)-f(u)=c(\bm{x})w$ (中值定理). 那么在 $\Sigma^{\prime}(\lambda)$ 内有
\begin{equation}\label{eq:radial-reflection-difference-equation}
\Delta w+c(\bm{x})w=0
\end{equation}
此外, 由引理的假设, $\Sigma^{\prime}(\lambda)$ 内有 $w\leq0$, 且 $w$ 在 $\Sigma^{\prime}(\lambda)$ 内不恒为零. 另外, $w$ 在 $T_\lambda\cap B$ 上为零, 这是 $\Sigma^{\prime}(\lambda)$ 边界的一部分. 由 Theorem~\ref{thm:strong-maximum-principle} 和 Lemma~\ref{lem:elliptic-boundary-point}, 得到在 $\Sigma^{\prime}(\lambda)$ 内 $w<0$, 且在 $T_\lambda\cap B$ 上 $w_1>0$. 由于 $T_\lambda\cap B$ 上 $w_1=-2u_1$, 引理成立.
\end{Proof}

\subsubsection{定理的证明}\label{subsec:radial-symmetry-theorem-proof}

用上一小节的两个引理证明, $u$ 关于平面 $x_1=0$ 对称, 且对 $\bm{x}\in B\cap\{x_1>0\}$ 有 $u_1<0$. 这显然蕴含定理, 因为由相同论证, $u$ 关于任意通过原点的平面都对称, 因而径向对称. 将证明以下结果.

\begin{Lemma}
\label{lem:moving-plane-comparison}
对任意 $\lambda\in(0,R)$, 有
\begin{equation}\label{eq:moving-plane-strict-comparison}
u_1(\bm{x})<0,\quad u(\bm{x})<u(\bm{x}^\lambda)\quad\forall\bm{x}\in\Sigma(\lambda).
\end{equation}

\end{Lemma}

由连续性, 在 $\Sigma(0)=B\cap\{x_1>0\}$ 内得到 $u(\bm{x})\leq u(\bm{x}^0)$. 将 $x_1$ 换成 $-x_1$, 重复相同论证, 得到在 $\Sigma(0)$ 内 $u(\bm{x})\geq u(\bm{x}^0)$. 因此 $u$ 关于平面 $x_1=0$ 对称, 定理成立.

\begin{Proof}
由 Lemma~\ref{lem:radial-boundary-directional-derivative}, 当 $\lambda$ 充分接近 $R$ 时, \eqref{eq:moving-plane-strict-comparison} 成立. 现在设 $\mu$ 为一个临界值, 使 \eqref{eq:moving-plane-strict-comparison} 对 $\lambda>\mu$ 成立, 但不能继续越过它. 由连续性, 得到
\begin{equation}\label{eq:moving-plane-critical-comparison}
u_1(\bm{x})<0,\quad u(\bm{x})\leq u(\bm{x}^\mu)\quad\forall\bm{x}\in\Sigma(\mu).
\end{equation}
需要证明 $\mu=0$. 反设 $\mu>0$. 对任意点 $\bm{x}_0\in\partial\Sigma(\mu)\setminus T_\mu$, 有 $\bm{x}_0^\mu\in B$, 因而 $0=u(\bm{x}_0)\ne u(\bm{x}_0^\mu)$. 所以 $u(\bm{x})$ 在 $\Sigma(\mu)$ 内不恒等于 $u(\bm{x}^\mu)$, 可以应用 Lemma~\ref{lem:radial-reflection-strictness}. 因而在 $\Sigma(\mu)$ 内 $u(\bm{x})<u(\bm{x}^\mu)$, 且在 $B\cap T_\mu$ 上 $u_1<0$. 所以 \eqref{eq:moving-plane-strict-comparison} 对 $\lambda=\mu$ 成立. 此外, 由 Lemma~\ref{lem:radial-boundary-directional-derivative}, $T_\mu\cap\partial B$ 上任意一点的某个邻域内都有 $u_1<0$. 因此 $T_\mu\cap\overline B$ 上每一点都具有使 $u_1<0$ 的邻域, 又因 $T_\mu\cap\overline B$ 是紧的, 必有对充分小的 $\epsilon$, 在 $\Sigma(\mu-\epsilon)$ 内 $u_1<0$.

由于假设了 \eqref{eq:moving-plane-strict-comparison} 不能继续越过 $\mu$, 必存在序列 $\lambda_j$ 及 $\bm{x}_j\in\Sigma(\lambda_j)$, 使 $\lambda_j\to\mu$, 且
\begin{equation}\label{eq:moving-plane-counterexample-sequence}
u(\bm{x}_j)\geq u(\bm{x}_j^{\lambda_j}).
\end{equation}
取子列后, 可假设 $\bm{x}_j$ 收敛; 极限 $\bm{x}$ 必位于 $\Sigma(\mu)$ 的闭包内. 在 \eqref{eq:moving-plane-counterexample-sequence} 中取极限, 得到 $u(\bm{x})\geq u(\bm{x}^\mu)$. 若 $\bm{x}\in\partial B$, 这不可能成立, 因为此时 $u(\bm{x})=0$, $u(\bm{x}^\mu)>0$. 由于 \eqref{eq:moving-plane-strict-comparison} 对 $\mu$ 成立, 必有 $\bm{x}\in T_\mu\cap B$; 因而 $\bm{x}^\mu=\bm{x}$. 另一方面, 连接 $\bm{x}_j$ 与 $\bm{x}_j^{\lambda_j}$ 的直线段位于 $B$ 内, 由 \eqref{eq:moving-plane-counterexample-sequence} 和中值定理, 其中必包含一点 $\bm{y}_j$, 使 $u_1(\bm{y}_j)\geq0$. 令 $j\to\infty$, 得到 $u_1(\bm{x})\geq0$. 矛盾.
\end{Proof}

\paragraph*{Problems}

\begin{enumerate}[label=\thesection.\arabic*.,ref=\thesection.\arabic*,start=17,series=book-problems]

\item\label{ex:nonnegative-nonlinearity-positive-solution} 证明, 若对每个 $u$ 都有 $f(u)\geq0$, 则 \eqref{eq:radial-symmetry-semilinear-equation}, \eqref{eq:radial-symmetry-zero-boundary} 的非平凡解自动为正. 再证明若 $f(0)\geq0$, 则 \eqref{eq:radial-symmetry-strict-decrease} 在 $r=R$ 处也成立.

\item\label{ex:radial-center-nonlinearity-positive} 证明, 若 \eqref{eq:radial-symmetry-semilinear-equation}, \eqref{eq:radial-symmetry-zero-boundary} 存在非负解 $u$, 且 $u$ 不恒为零, 则 $f(u(0))$ 必严格为正.

\item\label{ex:boundary-hessian-radial-identity} 设 $B$ 为以原点为球心的球. 设 $u$ 为 $C^2$ 函数, 在 $B$ 内 $u>0$, 且在 $\partial B$ 上 $u=0$. 设 $\bm{x}$ 为 $\partial B$ 上一点.

(a) 证明 $\Delta u(\bm{x})=u_{rr}(\bm{x})+\frac{n-1}{r}u_r(\bm{x})$, 其中 $r$ 为径向方向.

(b) 假设 $u_r(\bm{x})=0$. 证明 $\partial^2u/\partial x_i\partial x_j(\bm{x})=n_in_ju_{rr}(\bm{x})$. 这里 $\bm{n}$ 是 $\partial B$ 的单位法向. Hint: 先证明二阶导数构成的矩阵半正定.

\item\label{ex:radial-profile-semilinear-construction} 设 $u\in C^4(\overline B)$ 径向对称, 且当 $0<r\leq R$ 时满足 $u_r<0$, 并假设 $\Delta u(0)\ne0$. 证明存在 $C^1$ 函数 $f$, 使 $\Delta u+f(u)=0$.

\item\label{ex:annular-positive-solution-outward-monotonicity} 设 $u>0$ 为区域 $R_1<|\bm{x}|<R_2$ 内 \eqref{eq:radial-symmetry-semilinear-equation} 的一个 $C^2$ 解. 假设在外边界 $|\bm{x}|=R_2$ 上 $u=0$. 证明当 $(R_1+R_2)/2\leq|\bm{x}|<R_2$ 时, $u_r<0$.

\end{enumerate}
